\title{DeepSeekMath: Advancing the Boundaries of Mathematical Reasoning in Open Language Models}

\begin{abstract}
Mathematical reasoning represents a considerable challenge for language models due to its intricate and structured characteristics. In this work, we present DeepSeekMath~7B, which extends the pre-training of DeepSeek-Coder-Base-v1.5 7B using 120B math-related tokens derived from Common Crawl, combined with natural language and code datasets. DeepSeekMath~7B attains a remarkable score of 51.7\% on the challenging MATH benchmark without employing external tools or ensemble voting, nearing the capabilities of Gemini-Ultra and GPT-4. Utilizing self-consistency sampling with 64 samples from DeepSeekMath~7B yields a score of 60.9\% on MATH. The enhanced mathematical reasoning of DeepSeekMath~is credited to two main aspects: first, we exploit the vast potential of publicly accessible web data via a carefully designed data selection pipeline; second, we propose Group Relative Policy Optimization (GRPO), a PPO variant that not only improves mathematical reasoning but also optimizes PPO’s memory consumption.
\end{abstract}

\section{Introduction}

Large language models (LLMs) have transformed the methodology for mathematical reasoning in artificial intelligence, driving major progress on quantitative reasoning benchmarks \citep{MATH} and geometric reasoning challenges \citep{trinh2024solving}. Additionally, these models have demonstrated their utility in supporting humans to solve intricate mathematical problems \citep{tao}. Nevertheless, state-of-the-art models like GPT-4 \citep{gpt4} and Gemini-Ultra \citep{gemini} remain inaccessible to the public, and currently available open-source models lag significantly behind in terms of performance.

In this work, we present DeepSeekMath, a language model tailored to the mathematical domain that markedly surpasses the math abilities of existing open-source models and nearly matches the academic benchmark performance of GPT-4. To accomplish this, we develop the DeepSeekMath Corpus, a large-scale, high-quality pre-training dataset containing 120 billion math tokens. This corpus is derived from the Common Crawl (CC) via a fastText-based classifier \citep{joulin2016fasttext}. Initially, the classifier is trained using positive samples drawn from OpenWebMath \citep{openwebmath} and a varied set of other web pages as negative samples. Following this, we apply the classifier to identify more positive samples within the CC, which are subsequently refined through human annotation. The classifier is then retrained on this enhanced dataset to boost its accuracy. Evaluation demonstrates the high quality of the large corpus, with our DeepSeekMath-Base 7B model achieving 64.2\% accuracy on GSM8K \citep{gsm8k} and 36.2\% on the challenging MATH dataset \citep{MATH}, outperforming Minerva 540B \citep{minerva}. Moreover, since the DeepSeekMath Corpus is multilingual, we observe improvements on Chinese mathematical benchmarks \citep{wei2023cmath,agieval}. We consider our methodology for mathematical data curation as a foundation for the research community, with considerable potential for further enhancements.

DeepSeekMath-Base is initialized using DeepSeek-Coder-Base-v1.5 7B \citep{deepseek-coder}, as we find that beginning with a model trained on code is preferable to using a general LLM. Additionally, we note that mathematical training enhances the model's performance on MMLU \citep{mmlu} and BBH benchmarks \citep{bbh}, suggesting that it not only improves mathematical skills but also strengthens overall reasoning abilities.

Following pre-training, we perform mathematical instruction tuning on DeepSeekMath-Base employing chain-of-thought \citep{cot}, program-of-thought \citep{pot,pal}, and tool-integrated reasoning \citep{tora} datasets. The resulting DeepSeekMath-Instruct 7B model surpasses all other 7B models and is on par with 70B open-source instruction-tuned models.

Moreover, we propose the Group Relative Policy Optimization (GRPO), a variant of the Proximal Policy Optimization (PPO) \citep{schulman2017proximal} reinforcement learning (RL) algorithm. GRPO eliminates the need for a critic model by estimating the baseline using group scores, which considerably reduces the computational cost of training. Utilizing only a portion of English instruction tuning data, GRPO achieves significant gains over the strong DeepSeekMath-Instruct model, improving both in-domain tasks (GSM8K: 82.9\% $\rightarrow$ 88.2\%, MATH: 46.8\% $\rightarrow$ 51.7\%) and out-of-domain mathematical tasks (e.g., CMATH: 84.6\% $\rightarrow$ 88.8\%) during reinforcement learning. We also present a unified framework to interpret various approaches such as Rejection Sampling Fine-Tuning (RFT) \citep{yuan2023scaling}, Direct Preference Optimization (DPO) \citep{dpo}, PPO, and GRPO. Within this unified framework, we demonstrate that all these methods can be seen as either direct or simplified forms of RL. Extensive experiments, including comparisons of online versus offline training, outcome versus process supervision, and single-turn versus iterative RL, are conducted to thoroughly explore the core components of this framework. Finally, we discuss why our RL approach enhances the performance of instruction-tuned models and outline future directions for achieving more effective RL grounded in this unified framework.

\subsection{Contributions}
Our work presents scalable mathematical pre-training alongside an investigation and evaluation of reinforcement learning techniques.

\textbf{Math Pre-Training at Scale}

\begin{itemize}[topsep=0pt]
    \item We demonstrate strong evidence that the publicly available Common Crawl dataset holds valuable mathematical information. By employing a carefully crafted data selection process, we create the DeepSeekMath Corpus, a high-quality dataset comprising 120B tokens extracted from web pages filtered specifically for mathematical content. This dataset is nearly 7 times larger than the math-related web pages used by Minerva \citep{minerva} and 9 times larger than the recently introduced OpenWebMath \citep{openwebmath}.

    \item Our pre-trained base model, DeepSeekMath-Base 7B, attains performance comparable to Minerva 540B \citep{minerva}, suggesting that the parameter count alone does not determine mathematical reasoning ability. A smaller model trained on carefully curated data can also achieve strong results.

    \item We report insights from our math training experiments. Performing code training prior to math training enhances the models' capability to solve mathematical problems, both with and without the aid of tools. This partially addresses the ongoing debate: \textit{does code training enhance reasoning skills?} We argue that it does, at least in the context of mathematical reasoning.

    \item Although incorporating arXiv papers is a common practice, especially in math-centric studies, it does not yield significant improvements across the mathematical benchmarks used in our evaluations.
\end{itemize}

\textbf{Exploration and Analysis of Reinforcement Learning}

\begin{itemize}[topsep=0pt]
    \item We propose Group Relative Policy Optimization (GRPO), a reinforcement learning algorithm that is both efficient and effective. Unlike Proximal Policy Optimization (PPO), GRPO eliminates the need for a critic model by estimating the baseline using group scores, which considerably decreases the computational resources required for training.
    \item We show that GRPO substantially improves the performance of our instruction-tuned model, DeepSeekMath-Instruct, relying solely on instruction-tuning data. Additionally, we observe improvements in out-of-domain performance during the reinforcement learning phase.
    \item We offer a unified framework to comprehend various approaches, including RFT, DPO, PPO, and GRPO. We perform extensive experiments—such as online versus offline training, outcome versus process supervision, single-turn versus iterative reinforcement learning—to thoroughly explore the key components of this framework.
    \item Utilizing our unified framework, we investigate the underlying reasons for the success of reinforcement learning and outline several promising avenues for achieving more effective reinforcement learning of large language models (LLMs).
\end{itemize}

\subsection{Summary of Evaluations and Metrics}

\begin{itemize}[topsep=0pt]
    \item \textbf{English and Chinese Mathematical Reasoning}: We carry out thorough evaluations of our models on benchmarks in both English and Chinese, spanning mathematical problems from elementary to university level. The English benchmarks include GSM8K \citep{gsm8k}, MATH \citep{MATH}, SAT \citep{llemma}, OCW Courses \citep{minerva}, and MMLU-STEM \citep{mmlu}. Chinese benchmarks consist of MGSM-zh \citep{mgsm}, CMATH \citep{wei2023cmath}, Gaokao-MathCloze \citep{agieval}, and Gaokao-MathQA \citep{agieval}. We assess the models' capability to produce fully self-contained text solutions without relying on tools, as well as their ability to solve problems using Python.

    On the English datasets, DeepSeekMath-Base performs competitively against the closed-source Minerva 540B \citep{minerva}, outperforming all open-source base models (such as Mistral 7B \citep{mistral} and Llemma-34B \citep{llemma}) regardless of whether those models underwent math-specific pre-training, often by a notable margin. Significantly, DeepSeekMath-Base shows superior results on Chinese benchmarks, probably because our approach differs from prior work \citep{minerva,llemma} by including high-quality non-English math pre-training data rather than focusing exclusively on English. With the application of mathematical instruction tuning and reinforcement learning, the resulting DeepSeekMath-Instruct and DeepSeekMath-RL achieve strong outcomes, surpassing 50\% accuracy on the competitive MATH dataset—marking a first for open-source models.
\end{itemize}

\begin{itemize}
    \item \textbf{Formal Mathematics}: We assess DeepSeekMath-Base by performing the informal-to-formal theorem proving task from \citep{dsp_proof} on miniF2F \citep{minif2f}, with Isabelle \citep{isabelle} selected as the proof assistant. DeepSeekMath-Base shows strong few-shot autoformalization capabilities.
    \item \textbf{Natural Language Understanding, Reasoning, and Code}: To develop a thorough understanding of the models' overall comprehension, reasoning, and coding skills, we test DeepSeekMath-Base on the Massive Multitask Language Understanding (MMLU) benchmark \citep{mmlu}, which includes 57 multiple-choice tasks across various fields; BIG-Bench Hard (BBH) \citep{bbh}, containing 23 difficult tasks that primarily demand multi-step reasoning; and HumanEval \citep{codex} and MBPP \citep{mbpp}, both commonly used for evaluating code language models. Math pre-training enhances performance in both language understanding and reasoning.
\end{itemize}

\section{Math Pre-Training}

\subsection{Data Collection and Decontamination} This section describes the methodology for assembling the DeepSeekMath Corpus from Common Crawl. As illustrated in Figure \ref{fig:data_collect}, we introduce an iterative pipeline that systematically collects a large-scale mathematical dataset from Common Crawl, beginning with a seed corpus (for instance, a small yet high-quality math-related dataset). It is important to note that this method can be extended to other domains, such as programming.

Initially, we select OpenWebMath \citep{openwebmath}, a high-quality compilation of mathematical web texts, as our seed corpus. Leveraging this dataset, we train a fastText model \citep{joulin2016fasttext} to retrieve additional mathematical web pages similar to those in OpenWebMath. Concretely, we randomly pick 500,000 samples from the seed corpus as positive training examples, and 500,000 web pages from Common Crawl as negative samples. We use an open-source library\footnote{\url{https://fasttext.cc}} for training, setting the vector dimension to 256, the learning rate to 0.1, the maximum word n-gram length to 3, the minimum word frequency to 3, and training for 3 epochs. To reduce the size of the original Common Crawl, URL-based deduplication and near-deduplication techniques are applied, yielding 40B HTML web pages. We then retrieve mathematical web pages from this deduplicated Common Crawl using the fastText model. To remove low-quality mathematical content, the gathered pages are ranked according to the scores predicted by the fastText model, preserving only the highest-ranked ones. The amount of preserved data is evaluated by pre-training experiments on the top 40B, 80B, 120B, and 160B tokens. For the initial iteration, we retain the top 40B tokens.

Following the initial round of data gathering, many mathematical web pages remain uncollected, primarily due to the fastText model being trained on a set of positive samples that lack sufficient variety. To address this, we identify more mathematical web sources to augment the seed corpus, aiming to enhance the fastText model's performance. Concretely, we start by partitioning the entire Common Crawl dataset into distinct domains, where a domain is defined as web pages sharing the same base URL. For each domain, we compute the proportion of web pages collected during the first iteration. Domains with over 10\% of their pages collected are deemed math-related (for example, \url{mathoverflow.net}). Next, we manually label URLs within these domains that pertain to mathematical content (such as \url{mathoverflow.net/questions}). Web pages linked to these URLs that were not previously collected are incorporated into the seed corpus. This method allows us to accumulate additional positive examples, enabling the training of a more effective fastText model that can retrieve more mathematical data in future iterations. After four rounds of data collection, we obtain 35.5 million mathematical web pages amounting to 120 billion tokens. In the fourth iteration, we observe that nearly 98\% of the data had already been gathered by the third iteration, prompting us to halt further data collection.

To prevent contamination of benchmarks, we adopt the filtering strategy from \cite{deepseek-coder}, removing web pages containing questions or answers from English mathematical benchmarks such as GSM8K \citep{gsm8k} and MATH \citep{MATH}, as well as Chinese benchmarks like CMATH \citep{wei2023cmath} and AGIEval \citep{agieval}. The filtering process entails removing any text segment that contains a 10-gram string exactly matching any substring from the evaluation benchmarks from our math training corpus. For benchmark texts shorter than 10 grams but with at least 3 grams, we apply exact matching to exclude any contaminated web pages.

\subsection{Assessing the Quality of the DeepSeekMath~Corpus}

We conduct pre-training experiments to compare the DeepSeekMath~Corpus with recently released math-focused training corpora:

\begin{itemize}[topsep=0pt]
    \item \textbf{MathPile} \citep{mathpile}: a multi-source dataset (8.9B tokens) compiled from textbooks, Wikipedia, ProofWiki, CommonCrawl, StackExchange, and arXiv, with the majority (over 85\%) derived from arXiv;
    \item \textbf{OpenWebMath} \citep{openwebmath}: CommonCrawl data filtered to retain mathematical content, comprising 13.6B tokens;
    \item \textbf{Proof-Pile-2} \citep{llemma}: a mathematical dataset consisting of OpenWebMath, AlgebraicStack (10.3B tokens of mathematical code), and arXiv papers (28.0B tokens). For experiments on Proof-Pile-2, we follow \cite{llemma} by maintaining an arXiv:Web:Code ratio of 2:4:1.
\end{itemize}

\subsubsection{Training Setup}
\label{sec:quality-policy}
We conduct math training on a general pre-trained language model comprising 1.3B parameters, which follows the same architecture as the DeepSeek LLMs \citep{deepseek-llm}, referred to as DeepSeek-LLM 1.3B. Separate training is performed on each mathematical dataset for a total of 150B tokens. All experiments utilize the efficient and lightweight HAI-LLM \citep{haillm} training framework. Consistent with the training protocol of DeepSeek LLMs, we employ the AdamW optimizer \citep{adamW} with hyperparameters $\beta_1=0.9$, $\beta_2=0.95$, and $\mathrm{weight\_decay}=0.1$. We adopt a multi-step learning rate schedule where the learning rate increases to its maximum over 2,000 warmup steps, then decreases to 31.6\% of this peak after 80\% of training, and further drops to 10.0\% after 90\% of the training process. The maximum learning rate is set to 5.3e-4, and training is conducted with a batch size of 4 million tokens and a context length of 4K.

\subsubsection{Evaluation Outcomes}

\textbf{The DeepSeekMath~Corpus exhibits superior quality, supports multiple languages, and is the largest dataset available.}

\begin{itemize}[topsep=0pt]
    \item \textbf{Superior Quality}:
    We assess downstream performance using 8 mathematical benchmarks under few-shot chain-of-thought prompting \cite{cot}. As detailed in Table \ref{tab:corpora_comparison}, the model trained on the DeepSeekMath~Corpus consistently outperforms others. Figure \ref{fig:corpora_comparisons} further illustrates that at 50B tokens (equivalent to one complete epoch of Proof-Pile-2), the model trained on DeepSeekMath surpasses that trained on Proof-Pile-2, suggesting that the average quality of the DeepSeekMath Corpus is higher.
    \item \textbf{Multilingual}:
    The DeepSeekMath~Corpus includes multilingual content, with English and Chinese being the predominant languages. Table \ref{tab:corpora_comparison} shows that training on the DeepSeekMath~Corpus improves mathematical reasoning capabilities in both English and Chinese. In contrast, existing math corpora, which are mainly English-centric, tend to provide limited gains and can even impair performance in Chinese mathematical reasoning tasks.
    \item \textbf{Large-Scale}:
    The DeepSeekMath~Corpus is substantially larger than other existing mathematical datasets. As shown in Figure \ref{fig:corpora_comparisons}, DeepSeek-LLM 1.3B trained on the DeepSeekMath~Corpus exhibits a more pronounced learning curve with sustained improvement. By comparison, baseline corpora are considerably smaller, have been repeatedly cycled through during training, and their models quickly reach a performance plateau.
\end{itemize}

\subsection{Training and Evaluating DeepSeekMath-Base 7B}

In this part, we present DeepSeekMath-Base 7B, a foundational model with robust reasoning capabilities, particularly in mathematics. The model is initialized from DeepSeek-Coder-Base-v1.5 7B \citep{deepseek-coder} and trained on 500B tokens. The data composition includes: 56\% from the DeepSeekMath Corpus, 4\% from AlgebraicStack, 10\% from arXiv, 20\% from GitHub code, and the remaining 10\% consists of natural language data from Common Crawl in both English and Chinese. We largely follow the training configuration detailed in Section \ref{sec:quality-policy}, with the exceptions that the maximum learning rate is set to 4.2e-4 and the batch size is 10M tokens.

We perform an extensive evaluation of DeepSeekMath-Base 7B’s mathematical proficiency, concentrating on its capacity to generate self-contained mathematical solutions without external tools, solve math problems using tools, and carry out formal theorem proving. Additionally, we offer a broader overview of the base model’s abilities, including natural language comprehension, reasoning, and coding skills.

\paragraph{Mathematical Problem Solving with Step-by-Step Reasoning}

We assess DeepSeekMath-Base’s ability to solve math problems by employing few-shot chain-of-thought prompting \citep{cot} across eight benchmarks in both English and Chinese. These benchmarks cover quantitative reasoning tasks (e.g., GSM8K \citep{gsm8k}, MATH \citep{MATH}, and CMATH \citep{wei2023cmath}) as well as multiple-choice problems (e.g., MMLU-STEM \citep{mmlu} and Gaokao-MathQA \citep{agieval}), spanning a wide range of mathematical topics from basic to college-level complexity.

As illustrated in Table \ref{tab:base_math_cot}, DeepSeekMath-Base 7B achieves the highest performance across all eight benchmarks among open-source base models, which include the widely adopted general model Mistral 7B \citep{mistral} and the recently introduced Llemma 34B \citep{llemma} that was trained on math datasets from Proof-Pile-2 \citep{llemma}. Significantly, on the challenging MATH dataset, DeepSeekMath-Base surpasses all existing open-source base models by more than 10\% absolute and even exceeds Minerva 540B \citep{minerva}, a closed-source base model that is 77 times larger, built upon PaLM \citep{palm} and further trained on mathematical texts.

\paragraph{Mathematical Problem Solving with Tool Use} We assess program-aided mathematical reasoning on the GSM8K and MATH datasets employing few-shot program-of-thought prompting \citep{pot,pal}. The models are tasked with solving each problem by generating Python code, leveraging libraries such as \textit{math} and \textit{sympy} to perform complex calculations. The output from executing this code is then used as the answer. As detailed in Table \ref{tab:base_math_pot_proof}, DeepSeekMath-Base 7B outperforms the previous state-of-the-art, Llemma 34B.

\paragraph{Formal Mathematics}

Automating formal proofs is valuable for ensuring correctness and reliability in mathematical proofs while improving efficiency, an area that has garnered growing interest recently. We evaluate DeepSeekMath-Base 7B on the informal-to-formal proving task introduced by \citep{dsp_proof}, which requires generating a formal proof based on an informal statement, its formal equivalent, and an informal proof. This evaluation uses the miniF2F benchmark \citep{minif2f}, which focuses on formal Olympiad-level mathematics. For each problem, we generate formal proofs in Isabelle through few-shot prompting. Following \citep{dsp_proof}, models produce proof sketches, which are then completed by running the off-the-shelf automated prover Sledgehammer \citep{sledgehammer} to fill in missing details. As presented in Table \ref{tab:base_math_pot_proof}, DeepSeekMath-Base 7B shows strong capabilities in automatic proof formalization.

\paragraph{Natural Language Understanding, Reasoning, and Code}

We measure model abilities in natural language understanding via MMLU \citep{mmlu}, reasoning through BBH \citep{bbh}, and coding proficiency on HumanEval \citep{codex} and MBPP \citep{mbpp}. As reported in Table \ref{tab:understanding_reasoning_code}, DeepSeekMath-Base 7B achieves notable improvements over its predecessor, DeepSeek-Coder-Base-v1.5 \citep{deepseek-coder}, on MMLU and BBH, demonstrating the beneficial effect of math-focused training on language understanding and reasoning. Moreover, by incorporating code tokens during continual training, DeepSeekMath-Base 7B successfully preserves the coding performance of DeepSeek-Coder-Base-v1.5 on the two coding benchmarks. Overall, DeepSeekMath-Base 7B substantially outperforms the general-purpose model Mistral 7B \citep{mistral} across the three reasoning and coding tasks.

\section{Supervised Fine-Tuning}

\subsection{SFT Data Preparation}

We create a mathematical instruction-tuning dataset that includes English and Chinese problems from various mathematical disciplines and with different levels of difficulty: each problem is paired with solutions presented in chain-of-thought (CoT) \citep{cot}, program-of-thought (PoT) \citep{pot,pal}, and tool-assisted reasoning formats \citep{tora}. The dataset comprises a total of 776K training samples.

\begin{itemize}[topsep=0pt]
    \item \textbf{English mathematical datasets}: We annotate GSM8K and MATH problems with tool-integrated solutions, and use a portion of MathInstruct \citep{MathInstruct} along with the training portion of Lila-OOD \citep{lila}, where problems are solved using CoT or PoT methods. Our English dataset spans various branches of mathematics, such as algebra, probability, number theory, calculus, and geometry.
    \item \textbf{Chinese mathematical datasets}: We gather Chinese K-12 math problems covering 76 subcategories, including linear equations, with solutions annotated in both CoT and tool-assisted reasoning formats.
\end{itemize}

\subsection{Training and Evaluation of DeepSeekMath-Instruct 7B}

Here, we describe DeepSeekMath-Instruct 7B, which is fine-tuned for mathematical instruction tasks based on DeepSeekMath-Base. Training samples are concatenated randomly up to a maximum context length of 4K tokens. The model is trained for 500 steps, using a batch size of 256 and a fixed learning rate of 5e-5.

We assess the mathematical capabilities of models both without and with the utilization of tools, across four quantitative reasoning benchmarks in English and Chinese. Our model is compared with the leading models available at the time:

\begin{itemize}[topsep=0pt]
    \item \textbf{Closed-source models} include: (1) the GPT series, with GPT-4 \citep{gpt4} and GPT-4 Code Interpreter~\footnote{\url{https://openai.com/blog/chatgpt-plugins\#\#code-interpreter}} being the most advanced, (2) Gemini Ultra and Pro \citep{gemini}, (3) Inflection-2 \citep{inflection-2}, (4) Grok-1~\footnote{\url{https://x.ai/model-card}}, as well as recent models released by Chinese companies, including (5) Baichuan-3~\footnote{\url{https://www.baichuan-ai.com}}, and (6) the newest GLM-4~\footnote{\url{https://open.bigmodel.cn/dev/api\#glm-4}} from the GLM series \citep{glm}.

These models serve general purposes, with most having undergone multiple alignment procedures.
\item \textbf{Open-source models} encompass: general-purpose models such as (1) DeepSeek-LLM-Chat 67B \citep{deepseek-llm}, (2) Qwen 72B \citep{qwen}, (3) SeaLLM-v2 7B \citep{seallm}, and (4) ChatGLM3 6B \citep{chatglm3}, along with models enhanced for mathematics, including (5) InternLM2-Math 20B~\footnote{\url{https://github.com/InternLM/InternLM-Math}} which is based on InternLM2 and has been trained on math tasks followed by instruction tuning, (6) Math-Shepherd-Mistral 7B which applies PPO training \citep{schulman2017proximal} to Mistral 7B \citep{mistral} using a process-supervised reward model, (7) the WizardMath series \citep{wizardmath} that improves mathematical reasoning capabilities in Mistral 7B and Llama-2 70B \citep{llama2} through evolve-instruct (a variant of instruction tuning using AI-evolved instructions) and PPO training, relying primarily on training problems from GSM8K and MATH, (8) MetaMath 70B \citep{metamath}, which fine-tunes Llama-2 70B on an augmented GSM8K and MATH dataset, (9) ToRA 34B \cite{tora}, a CodeLlama 34B fine-tuned for tool-integrated mathematical reasoning, and (10) MAmmoTH 70B \citep{MathInstruct}, which instruction-tunes Llama-2 70B on MathInstruct.
\end{itemize}

As indicated in Table \ref{tab:sft_rl_math}, when evaluated without tool usage, DeepSeekMath-Instruct 7B shows robust step-by-step reasoning performance. Remarkably, on the competition-level MATH dataset, this model outperforms all open-source models and most proprietary ones (such as Inflection-2 and Gemini Pro) by an absolute margin of at least 9\%. This advantage holds even against substantially larger models (e.g., Qwen 72B) and those specifically enhanced with math-focused reinforcement learning (e.g., WizardMath-v1.1 7B). Although DeepSeekMath-Instruct competes closely with Chinese proprietary models GLM-4 and Baichuan-3 on MATH, it still falls short of GPT-4 and Gemini Ultra.

In the evaluation setting permitting both natural language reasoning and program-based tool use for solving problems, DeepSeekMath-Instruct 7B achieves nearly 60\% accuracy on MATH, outperforming all existing open-source models. On other benchmarks, our model performs comparably to DeepSeek-LLM-Chat 67B, the previous state-of-the-art that is ten times larger.
\section{Reinforcement Learning}

\subsection{Group Relative Policy Optimization} Reinforcement learning (RL) has demonstrated effectiveness in further enhancing the mathematical reasoning abilities of LLMs after Supervised Fine-Tuning (SFT) \citep{wang2023math,wizardmath}. In this section, we present our efficient and effective RL algorithm, Group Relative Policy Optimization (GRPO).

\subsubsection{From PPO to GRPO} Proximal Policy Optimization (PPO) \citep{schulman2017proximal} is an actor-critic RL method commonly employed during RL fine-tuning of LLMs \citep{ouyang2022training}. Specifically, it optimizes LLMs by maximizing the following surrogate objective:

\begin{equation}
\footnotesize
    \mathcal{J}_{PPO}(\theta) = \mathbb{E}{[q \sim P(Q), o \sim \pi_{\theta_{old}}(O|q)]} \frac{1}{|o|} \sum_{t=1}^{|o|} \min \left[ \frac{\pi_\theta(o_{t} | q, o_{<t})}{\pi_{\theta_{old}}(o_{t} | q, o_{<t})} A_{t}, \text{clip} \left( \frac{\pi_\theta(o_{t} | q, o_{<t})}{\pi_{\theta_{old}}(o_{t} | q, o_{<t})}, 1 - \varepsilon, 1 + \varepsilon \right)  A_{t} \right] ,
\end{equation}

Here, $\pi_{\theta}$ and $\pi_{\theta_{old}}$ represent the current and previous policy models, respectively, while $q$ and $o$ denote questions and outputs sampled from the question dataset and the old policy $\pi_{\theta_{old}}$. The parameter $\varepsilon$ is a clipping-related hyper-parameter introduced in PPO to stabilize training. The advantage $A_t$ is computed using Generalized Advantage Estimation (GAE) \citep{gae}, relying on the rewards $\{r_{\ge t}\}$ and a learned value function $V_{\psi}$. Consequently, PPO requires training a value function alongside the policy model. To prevent excessive optimization of the reward model, the standard practice is to include a per-token KL penalty relative to a reference model within the reward at each token \citep{ouyang2022training}, as expressed by

\begin{equation}
    r_{t} = r_\varphi(q, o_{\le t}) - \beta \log\frac{\pi_{\theta}(o_{t}|q, o_{<t})}{\pi_{ref}(o_{t}|q, o_{<t})},
\label{eq:PPO-reward}
\end{equation}

where $r_\varphi$ denotes the reward model, $\pi_{ref}$ is the reference model, typically the initial SFT model, and $\beta$ controls the KL penalty strength.

Since the value function in PPO is usually a model of similar size to the policy model, it imposes considerable memory and computational costs. Moreover, during RL training, the value function serves as a baseline in advantage estimation to reduce variance. However, in the context of LLMs, the reward model commonly assigns a score only to the final token, complicating training a value function that accurately estimates token-level values. To overcome this, as illustrated in Figure \ref{fig:grpo}, we introduce Group Relative Policy Optimization (GRPO), which eliminates the need for a separate value function approximation as required in PPO. Instead, it utilizes the average reward of multiple outputs generated from the same question as a baseline. Specifically, for each question $q$, GRPO samples a set of outputs $\{o_1, o_2, \cdots, o_G\}$ from the old policy $\pi_{\theta_{old}}$ and optimizes the policy by maximizing

\begin{equation}
\footnotesize
\begin{split}
    \mathcal{J}_{GRPO}(\theta) &= \mathbb{E}{[q \sim P(Q), \{o_i\}_{i=1}^G \sim \pi_{\theta_{old}}(O|q)]}  \\
    & \frac{1}{G}\sum_{i=1}^G\frac{1}{|o_i|} \sum_{t=1}^{|o_i|} \left\{ \min \left[ \frac{\pi_\theta(o_{i,t} | q, o_{i,<t})}{\pi_{\theta_{old}}(o_{i,t} | q, o_{i,<t})} \hat{A}_{i,t}, \text{clip} \left( \frac{\pi_\theta(o_{i,t} | q, o_{i,<t})}{\pi_{\theta_{old}}(o_{i,t} | q, o_{i,<t})}, 1 - \varepsilon, 1 + \varepsilon \right)  \hat{A}_{i,t} \right] - \beta \mathbb{D}_{KL}\left[\pi_{\theta} || \pi_{ref}\right]\right\} ,
\end{split}
\label{eq:GRPO-obj}
\end{equation}

where $\varepsilon$ and $\beta$ are hyper-parameters, and $\hat{A}_{i,t}$ represents the advantage computed solely from the relative rewards among outputs within each group, as detailed in subsequent sections. The group-relative method for advantage calculation in GRPO aligns well with the comparative nature of reward models, which are usually trained on datasets comprising comparisons of outputs for the same question. Notably, unlike PPO where the KL penalty is added inside the reward, GRPO applies regularization by directly incorporating the KL divergence between the trained policy and the reference policy into the loss function, simplifying the computation of $\hat{A}_{i,t}$. Additionally, in contrast to the KL penalty term in (\ref{eq:PPO-reward}), the KL divergence here is estimated via the following unbiased estimator \citep{kl_approx}:

\begin{equation}
\small
    \mathbb{D}_{KL}\left[\pi_{\theta} || \pi_{ref}\right] =

\begin{algorithm}[t]
  \small
  \caption{Iterative Group Relative Policy Optimization}
  \textbf{Input} initial policy model $\pi_{\theta_{\text{init}}}$; reward models $r_\varphi$; task prompts $\mathcal{D}$; 
  hyperparameters $\varepsilon$, $\beta$, $\mu$
  \begin{algorithmic}[1]
    \State Initialize policy model $\pi_\theta \leftarrow \pi_{\theta_{\text{init}}}$
    \For{iteration = 1, \dots, I}
       \State Set reference model $\pi_{ref} \leftarrow \pi_{\theta}$
      \For{step = 1, \dots, M}
      \State Draw a batch $\mathcal{D}_b$ from $\mathcal{D}$
      \State Save current policy as old policy $\pi_{\theta_{old}} \leftarrow \pi_{\theta}$ 
      \State Generate $G$ outputs $\{o_i\}_{i=1}^G \sim \pi_{\theta_{old}} (\cdot \mid q) $ for each question $q \in \mathcal{D}_b$
      \State Evaluate rewards $\{r_i\}_{i=1}^{G}$ for each output $o_i$ using $r_{\varphi}$ 
      \State Calculate $\hat{A}_{i,t}$ for the $t$-th token in $o_i$ via group relative advantage estimation.
      \For{GRPO iteration = 1, \dots, $\mu$}
        \State Update policy $\pi_{\theta}$ by maximizing the GRPO objective (Equation \ref{eq:GC-GRPO})
      \EndFor
    \EndFor \State Update reward model $r_\varphi$ continuously using a replay-based training scheme. 
    \EndFor 
  \end{algorithmic}
  \textbf{Output} $\pi_\theta$
  \label{alg:iter-grpo}
\end{algorithm}

\subsubsection{Outcome Supervision RL with GRPO} Formally, for each question $q$, a group of outputs $\{o_1, o_2, \ldots, o_G\}$ are sampled from the old policy $\pi_{\theta_{old}}$. A reward model assigns scores to these outputs, producing a reward set $\mathbf{r}=\{r_1, r_2, \ldots, r_G\}$. These rewards are then normalized by subtracting the mean of the group and dividing by the group standard deviation. Outcome supervision provides the normalized reward only at the completion of each output $o_i$ and assigns the advantage values $\hat{A}_{i, t}$ for every token in the output to be equal to this normalized reward, i.e., $\hat{A}_{i, t} = \widetilde{r}_i = \frac{r_i- {\rm mean}(\mathbf{r})}{{\rm std}(\mathbf{r})}$, before optimizing the policy according to the objective defined in equation (\ref{eq:GRPO-obj}).

\subsubsection{Process Supervision RL with GRPO} Since outcome supervision only delivers a reward at the end of each generated output, it might be inadequate or inefficient for guiding the policy in complicated mathematical tasks. Inspired by \cite{wang2023math}, we also examine process supervision, which assigns rewards at the conclusion of each reasoning step. Specifically, for a question $q$ and its $G$ sampled outputs $\{o_1, o_2, \ldots, o_G\}$, a process reward model scores every step within these outputs, resulting in reward sets: $\mathbf{R} = \{ \{r_1^{index(1)},\ldots,r_1^{index(K_1)}\}, \ldots,  \{r_G^{index(1)},\ldots,r_G^{index(K_G)}\} \}$, where $index(j)$ denotes the end token index of the $j$-th step, and $K_i$ is the number of steps in output $i$. These rewards are normalized across all steps by subtracting the mean and dividing by the standard deviation, i.e., $\widetilde{r}_i^{index(j)} = \frac{r_i^{index(j)} - {\rm mean}(\mathbf{R})}{{\rm std}(\mathbf{R})}$. The advantage for each token is then computed as the sum of the normalized rewards for all subsequent steps from token $t$ onward, formally $\hat{A}_{i, t} = \sum_{index(j) \ge t} \widetilde{r

\subsubsection{Iterative RL with GRPO} As the reinforcement learning training advances, the previous reward model may become inadequate for guiding the updated policy model. Consequently, we investigate iterative RL employing GRPO. As detailed in Algorithm \ref{alg:iter-grpo}, during iterative GRPO, new training datasets for the reward model are generated based on samples produced by the policy model. The existing reward model is then continuously refined using a replay strategy that incorporates 10\% of past data. Subsequently, the policy model is set as the reference model, and it is iteratively trained utilizing the updated reward model.

\subsection{Training and Evaluating DeepSeekMath-RL}

Our RL experiments utilize DeepSeekMath-Instruct 7B. The RL training data consist of chain-of-thought formatted questions related to GSM8K and MATH drawn from the SFT dataset, totaling approximately 144K questions. We exclude other SFT questions to specifically assess the effect of RL on benchmarks with limited data during the RL phase. The construction of the reward model training set follows the method of \citep{wang2023math}. The initial reward model is trained based on DeepSeekMath-Base 7B with a learning rate of 2e-5. For GRPO, the policy model’s learning rate is set to 1e-6, and the KL divergence coefficient is fixed at 0.04. For each question, we sample $64$ outputs. The maximum sequence length is set to 1024, with a training batch size of 1024. The policy model undergoes only one update after each exploration stage. DeepSeekMath-RL 7B is evaluated on benchmarks similarly to DeepSeekMath-Instruct 7B. For DeepSeekMath-RL 7B, tasks such as GSM8K and MATH that include chain-of-thought reasoning are treated as in-domain, while all other benchmarks are considered out-of-domain.

Table \ref{tab:sft_rl_math} presents the results of open- and closed-source models evaluated on English and Chinese benchmarks, employing both chain-of-thought and tool-integrated reasoning approaches. Our observations include: 1) DeepSeekMath-RL 7B achieves accuracy rates of 88.2\% on GSM8K and 51.7\% on MATH through chain-of-thought reasoning. This performance exceeds that of all open-source models within the 7B to 70B parameter range, as well as most closed-source counterparts. 2) Importantly, DeepSeekMath-RL 7B was trained solely on chain-of-thought-format instruction tuning data from GSM8K and MATH, initialized from DeepSeekMath-Instruct 7B. Despite the limited training dataset, it surpasses DeepSeekMath-Instruct 7B on every evaluation metric, demonstrating the effectiveness of reinforcement learning.

\section{Discussion}
In this section, we present insights derived from our pre-training and reinforcement learning experiments.

\subsection{Insights Gained from Pre-Training}

We begin by discussing our pre-training experiences. Unless otherwise noted, training configurations follow those described in Section \ref{sec:quality-policy}. Notably, when mentioning the DeepSeekMath Corpus here, we refer to the 89B-token dataset obtained from the second round of data collection.

\subsubsection{Code Training Enhances Mathematical Reasoning}

A widely discussed but unconfirmed idea is that training on code improves reasoning capabilities. We aim to provide partial evidence supporting this notion, specifically within the mathematical reasoning domain: code training enhances models' proficiency in mathematical reasoning both with and without the assistance of tools.

To investigate the impact of code training on mathematical reasoning, we conducted experiments using both two-stage and single-stage training frameworks:

\textbf{Two-Stage Training}

\begin{itemize}[topsep=0pt]
    \item \textbf{Code Training for 400B Tokens $\rightarrow$ Math Training for 150B Tokens}: DeepSeek-LLM 1.3B is initially trained on 400B code tokens, followed by 150B math tokens;
    \item \textbf{General Training for 400B Tokens $\rightarrow$ Math Training for 150B Tokens}: As a baseline comparison, we replace code tokens with general tokens (drawn from a large-scale general corpus compiled by DeepSeek-AI) during the first training phase, aiming to explore whether code tokens provide advantages over general tokens for enhancing mathematical reasoning.
\end{itemize}

\textbf{One-Stage Training}

\begin{itemize}[topsep=0pt]
    \item \textbf{Math Training for 150B Tokens}: DeepSeek-LLM 1.3B is trained solely on 150B math tokens;
    \item \textbf{Training on a mixture of 400B Code Tokens and 150B Math Tokens}: Since math training following code training tends to reduce coding performance, we examine if training on a combined dataset of code and math tokens in a single stage can still boost mathematical reasoning while preventing catastrophic forgetting.
\end{itemize}

\paragraph{Results}
Tables \ref{tab:code-math} and \ref{tab:code-math-general-eval} present the downstream performance across various training configurations.

Training on code improves program-assisted mathematical reasoning in both two-stage and one-stage training scenarios. As indicated in Table \ref{tab:code-math}, within the two-stage training framework, code training alone substantially elevates the capacity to solve GSM8K and MATH tasks using Python, and subsequent math training further enhances results. Notably, in the one-stage training scenario, combining code and math tokens effectively reduces catastrophic forgetting associated with two-stage training and benefits both coding (Table \ref{tab:code-math-general-eval}) and program-guided mathematical reasoning (Table \ref{tab:code-math}).

Code training also enhances mathematical reasoning without the aid of tools. In the two-stage setting, the initial code training phase already produces moderate gains. It also accelerates the effectiveness of the later math training stage, ultimately achieving the highest performance. Conversely, mixing code and math tokens for one-stage training diminishes mathematical reasoning without tool use. This may be because DeepSeek-LLM 1.3B, given its relatively small model size, lacks sufficient capacity to simultaneously integrate both code and mathematical data.

\subsubsection{ArXiv Papers Appear Ineffective in Enhancing Mathematical Reasoning}

ArXiv papers are frequently used as part of the math pre-training datasets \citep{minerva,gpt-f,llemma,mathpile}. Nonetheless, a thorough investigation into their influence on mathematical reasoning has not been widely performed. Interestingly, contrary to expectations, our experiments suggest that arXiv papers do not effectively boost mathematical reasoning capabilities. We conduct experiments with models of various sizes, including DeepSeek-LLM 1.3B and DeepSeek-Coder-Base-v1.5 7B \citep{deepseek-coder}, employing arXiv datasets processed through different pipelines:

\begin{itemize}[topsep=0pt]
    \item \textbf{MathPile} \citep{mathpile}: an 8.9 billion-token dataset curated with heuristic cleaning and filtering, composed of over 85\% scientific arXiv papers;
    \item \textbf{ArXiv-RedPajama} \citep{redpajama}: a collection of all arXiv LaTeX files with preambles, comments, macros, and bibliographies removed, amounting to 28.0 billion tokens.
\end{itemize}

In our tests, we independently train DeepSeek-LLM 1.3B for 150 billion tokens and DeepSeek-Coder-Base-v1.5 7B for 40 billion tokens on each arXiv corpus. The results suggest that arXiv papers do not effectively improve mathematical reasoning. When trained solely on arXiv data, both models show no significant progress or even regression across a range of mathematical benchmarks with varying difficulty levels utilized in this study. These benchmarks encompass quantitative reasoning datasets such as GSM8K and MATH (see Table \ref{tab:arxiv-cot}), multiple-choice tasks like MMLU-STEM (Table \ref{tab:arxiv-cot}), and formal mathematics benchmarks such as miniF2F (Table \ref{tab:arxiv_atp}).

However, this conclusion should be interpreted cautiously due to certain limitations. We have not yet examined:

\begin{itemize}[topsep=0pt]
    \item The influence of arXiv tokens on particular math-related tasks excluded from this study, for example, informalization of theorems, which involves converting formal statements or proofs into informal versions;
    \item The effect of arXiv tokens when integrated with other data types;
    \item Whether the advantages of arXiv papers emerge when applied to models of larger scale.
\end{itemize}

Therefore, additional investigation is warranted, which we defer to future work.

\subsection{Insights of Reinforcement Learning}
\subsubsection{Towards a Unified Paradigm} In this section, we present a unified framework for analyzing various training approaches, including SFT, RFT, DPO, PPO, and GRPO, and conduct experiments to investigate components of this unified paradigm. Generally, the gradient with respect to the parameter $\theta$ for a training method can be expressed as:

\begin{equation}
    \nabla_{\theta}\mathcal{J}_{\textcolor{red}{\mathcal{A}}}(\theta) = \mathbb{E}[\underbrace{(q,o) \sim \textcolor{red}{\mathcal{D}}}_{Data \ Source}]\left( \frac{1}{|o|} \sum_{t=1}^{|o|}  \underbrace{GC_{{\mathcal{A}}}(q, o, t, \textcolor{red}{\pi_{{rf}}})}_{Gradient \ Coefficient}  \nabla_{\theta}\log \pi_{\theta}(o_t | q, o_{<t})\right).
\label{eq:objective}
\end{equation}

There are three fundamental elements: 1) \textit{Data Source $\mathcal{D}$}, which defines the dataset used for training; 2) \textit{Reward Function $\pi_{{rf}}$}, providing the reward signal during training; 3) \textit{Algorithm $\mathcal{A}$}, which utilizes the training data and reward signals to compute the gradient coefficient $GC$, indicating the scale of penalty or reward applied to the data. We examine several representative approaches framed within this unified structure:

\begin{itemize}
    \item  \textbf{Supervised Fine-tuning (SFT)}: SFT adapts the pretrained model by fine-tuning it on human-curated SFT datasets.
    \item \textbf{Rejection Sampling Fine-tuning (RFT)}: RFT performs additional fine-tuning on the SFT model using outputs sampled from the SFT model’s responses to SFT questions, filtering these outputs based on the accuracy of their answers.
    \item \textbf{Direct Preference Optimization (DPO)}: DPO further improves the SFT model by fine-tuning it with augmented outputs sampled from the SFT model, employing a pairwise DPO loss function.
    \item \textbf{Online Rejection Sampling Fine-tuning (Online RFT)}: Unlike RFT, Online RFT starts from the SFT model as the initial policy and fine-tunes it with augmented outputs generated from the current policy model in real time.
    \item \textbf{PPO/GRPO}: PPO/GRPO also initializes the policy from the SFT model and strengthens it using outputs sampled from the current policy model.
\end{itemize}

We provide a summary of the elements of these methods in Table \ref{tab:data_policy}. For a more comprehensive derivation process, please consult Appendix \ref{app:analysis-rl}.

\paragraph{Observation regarding Data Source} We categorize the data source into two types: online sampling and offline sampling. Online sampling refers to training data obtained from exploration results generated by the current training policy model, whereas offline sampling refers to training data derived from the sampling outcomes of the initial SFT model. RFT and DPO utilize offline sampling, while Online RFT and GRPO employ online sampling.

As illustrated in Figure \ref{fig:rl-analysis}, we observe that Online RFT significantly outperforms RFT across two benchmarks. Notably, Online RFT performs similarly to RFT in the early training stages but achieves a clear advantage in later stages, indicating the benefits of online training. This observation is intuitive because, at the beginning, the actor and the SFT model are closely aligned, resulting in sampled data that differ only slightly. However, in later stages, data sampled from the actor show more pronounced differences, making real-time data sampling more beneficial.

\paragraph{Observation regarding Gradient Coefficient} The algorithm transforms the reward signal into a gradient coefficient to update the model parameters. In our experiments, we classify the reward function into two types: `Rule' and `Model'. `Rule' involves assessing the quality of a response based on the correctness of the answer, while `Model' means that we train a reward model to assign scores to each response. The reward model’s training data is derived from rule-based judgments. Equations \ref{eq:GC-RFT} and \ref{eq:GC-GRPO} emphasize a key distinction between GRPO and Online RFT: GRPO uniquely modifies its gradient coefficient according to the reward value from the reward model, enabling differential reinforcement and penalization of responses based on their varying reward magnitudes. Conversely, Online RFT lacks this capability; it does not penalize incorrect responses and applies uniform reinforcement to all responses with correct answers at the same intensity level.

As illustrated in Figure \ref{fig:rl-analysis}, GRPO outperforms online RFT, underscoring the effectiveness of adjusting the coefficients of positive and negative gradients. Moreover, GRPO+PS exhibits better performance than GRPO+OS, demonstrating the advantages of employing fine-grained, step-aware gradient coefficients. Additionally, we investigate iterative RL, performing two rounds of iteration in our experiments. As depicted in Figure \ref{fig:iter-rl}, iterative RL substantially enhances performance, particularly during the first iteration.

\subsubsection{Why Does RL Work?} In this study, we apply reinforcement learning to a subset of instruction tuning data, resulting in a significant performance improvement over the instruction-tuned model. To further clarify why reinforcement learning is effective, we assess the Pass@K and Maj@K accuracy of both the Instruct and RL models on two benchmark datasets. Figure \ref{fig:maj-pass} shows that RL boosts the Maj@K performance but does not improve Pass@K. These results suggest that RL improves the model’s overall robustness in output distribution, implying that \textbf{the enhancement is likely due to increasing the likelihood of the correct response appearing within the TopK rather than fundamentally improving the model’s capabilities.} Similarly, \citep{wang2023making} identified a \textbf{misalignment issue} in reasoning tasks for the SFT model, revealing that reasoning performance can be enhanced through a series of preference alignment approaches \citep{yuan2023rrhf,song2023preference,wang2023making}.

\subsubsection{How Can More Effective RL Be Achieved?}
\label{sec:effec_rl}
We have shown that RL is highly effective for mathematical reasoning tasks. Furthermore, we propose a unified framework to interpret various representative training methods, conceptualizing them as either direct or simplified RL approaches. As outlined in Equation \ref{eq:objective}, three fundamental components are involved: Data Source, Algorithm, and Reward Function. We suggest several promising future directions related to these components.

\paragraph{Data Source} The data source serves as the foundational input for all training methodologies. Within RL, we specifically define the data source as unlabeled questions accompanied by outputs sampled from the policy model. In this work, we exclusively utilize questions from the instruction tuning phase and apply a simple nucleus sampling technique to generate outputs. We believe this might explain why our RL framework only enhances Maj@K performance. Going forward, we plan to extend our RL approach to out-of-distribution question prompts, combined with \textbf{advanced sampling (decoding) strategies} such as those leveraging tree-search algorithms \citep{yao2023tree}. Additionally, \textbf{efficient inference methods} \citep{xia-etal-2023-speculative,leviathan2023fast,kwon2023efficient,xia2024unlocking}, which directly affect the exploration efficiency of policy models, will also be crucial.

\paragraph{Algorithms} Algorithms utilize the data and reward signal to determine the gradient coefficient used for updating the model parameters. According to Equation \ref{eq:objective}, all current methods to some degree fully \textbf{TRUST} the reward function's signal to either increase or decrease the conditional probability of a specific token. However, it is not always possible to guarantee that the reward signal is consistently reliable, particularly in highly complex tasks. For instance, even the PRM800K datasets \citep{lightman2023let}, which have been meticulously labeled by skilled annotators, still contain roughly 20\% of incorrect annotations\footnote{\url{https://github.com/openai/prm800k/issues/12\#issuecomment-1728491852}}. To address this, we aim to investigate reinforcement learning algorithms that are robust in the presence of noisy reward signals. We believe that such \textbf{WEAK-TO-STRONG} \citep{burns2023weak} alignment approaches will fundamentally transform learning algorithms.

\paragraph{Reward Function} The reward function serves as the origin of the training signal. In reinforcement learning, the reward function is typically represented by a neural reward model. We identify three crucial directions for advancing reward models: 1) \textbf{How to improve the generalization capability of the reward model.} The reward model must generalize effectively to handle out-of-distribution queries and complex decoding results; otherwise, reinforcement learning might only stabilize the output distribution of LLMs without enhancing their core abilities; 2) \textbf{How to incorporate the uncertainty of the reward model.} Capturing uncertainty could potentially bridge weak reward models and weak-to-strong learning algorithms; 3) \textbf{How to efficiently develop high-quality process reward models} that can provide detailed, fine-grained training signals for reasoning processes \citep{lightman2023let,wang2023math}.

\section{Conclusion, Limitation, and Future Work} We introduce DeepSeekMath, which surpasses all open-source models on the competition-level MATH benchmark and nears the performance of proprietary models. DeepSeekMath~is initially based on DeepSeek-Coder-v1.5 7B and is continually trained on 500B tokens, with a substantial portion of the data comprising 120B mathematical tokens extracted from Common Crawl. Our comprehensive ablation study indicates that web pages provide considerable promise for acquiring high-quality mathematical data, whereas arXiv appears to be less advantageous than initially anticipated. We propose Group Relative Policy Optimization (GRPO), a variant of Proximal Policy Optimization (PPO), which significantly enhances mathematical reasoning capabilities while reducing memory usage. Experimental results demonstrate that GRPO remains effective even when DeepSeekMath-Instruct 7B has achieved high benchmark scores. Furthermore, we offer a unified framework to interpret a range of methods and outline several promising directions for more efficient reinforcement learning.

Despite DeepSeekMath~achieving strong results on quantitative reasoning benchmarks, its performance in geometry and theorem-proving tasks is comparatively weaker than that of proprietary models. For example, in our preliminary testing, the model struggles with problems involving triangles and ellipses, suggesting a possible data selection bias during pre-training and fine-tuning. Additionally, due to limitations in model size, DeepSeekMath~underperforms GPT-4 in few-shot learning scenarios. While GPT-4 demonstrates improved performance when provided with few-shot inputs, DeepSeekMath~exhibits similar performance across both zero-shot and few-shot evaluations. Moving forward, we plan to enhance our engineered data selection pipeline to build a higher-quality pre-training corpus. We will also investigate the potential avenues discussed in Section \ref{sec:effec_rl} for more effective reinforcement learning in large language models.

\end{document}